\section{Conclusão}
\label{sec:conclusao}

A avaliação transformou percepções sobre o SIGAA/UnB em evidências comparadas a critérios fixados antes da coleta. A resposta à pergunta da Fase 1 é que \textbf{o SIGAA entrega o que o estudante precisa, mas a um custo alto de tempo, esforço e satisfação, e não foi feito para o celular, o aparelho que todos os participantes usam}. O núcleo funcional é sólido: as funções das seis tarefas existem, os dados acadêmicos estão corretos, 83,3\% das tentativas terminam sem ajuda e as tarefas funcionam em todos os navegadores de computador. O custo está na interação, com tempo de 4,39 vezes o do especialista e SUS de 32,5, metade da média de referência; a acessibilidade exclui quem não usa mouse; e nenhuma página se adapta a telas pequenas. A matriz multiplataforma mostrou que esse resultado não depende do sistema, do aparelho nem do navegador, e que a página de login da UnB, no mesmo percurso, já é responsiva, o que indica que a correção é viável no próprio ambiente da universidade.

\textbf{Melhorias recomendadas.} Os resultados apoiam, em ordem de prioridade, as seguintes decisões da STI/UnB, do DEG e das coordenações:
\begin{enumerate}
	\item declarar a \textit{viewport} e substituir a largura fixa de 980~px por um \textit{layout} fluido no portal do discente (M14, matriz multiplataforma e IQ4);
	\item tornar o menu principal operável pelo teclado, pelo leitor de tela e pelo toque (M10 e M14);
	\item corrigir o link do calendário no painel do portal e destacar a carga horária pendente, as tarefas de pior índice (T6 e T4);
	\item mostrar o horário por extenso ao lado do código (M11) e corrigir os erros do WAVE, concentrados na Turma Virtual e no portal (M9);
	\item avisar a expiração da sessão (M20) e reduzir o peso da Turma Virtual, que passa de 19~s no celular em 4G;
	\item até que as correções cheguem, orientar os calouros sobre o calendário, a integralização e o código de horário; e, a médio prazo, oferecer o planejamento do fluxo do curso, que leva 60\% dos participantes a ferramentas externas (M3).
\end{enumerate}

\textbf{Limitações.} A amostra do teste com usuários tem cinco participantes recrutados por conveniência na FCTE, com dois participantes no grupo do começo do curso; M14 e a matriz multiplataforma foram feitas por emulação; M15 teve um avaliador; os tempos de desempenho vêm de um único ponto de rede e de horários fora do pico; a matrícula e o trancamento não foram avaliados, porque alterariam dados reais; e a equipe não teve acesso ao código-fonte.

\textbf{Trabalhos futuros.} Repetir o teste com usuários no protótipo, para medir o efeito da melhoria no tempo, na conclusão e na satisfação; confirmar a matriz em aparelhos reais de cada plataforma; ampliar a amostra para outros \textit{campi} e para estudantes com deficiência que usam leitor de tela; e medir o desempenho nos períodos de matrícula, quando o servidor está sob carga.
